\documentclass{ip-quiz}

\quizlevel{n2}
\quiznum{1}
\quiztitle{Funciones y strings}

\begin{document}

\makeheader
\maketitle

\begin{sectionbox}{Enunciado}
Escriba una función que recibe como parémetro una cadena de caracteres alfanumérica y retorne un histograma, representado como un diccionario, de las apariciones de cada dígito en la cadena.
\end{sectionbox}

\begin{sectionbox}{Solución}
\begin{minted}[escapeinside=||]{python}
def histograma_digitos(cadena: str) -> dict:
    histograma = {}
    i = 0
    while i < len(cadena):
        if |\tms{isdig}|cadena[i].isdigit():|\tme{isdig}|
            histograma[cadena[i]] = histograma.|\tms{get}|get(cadena[i], 0)|\tme{get}| + 1
        i += 1
    return histograma
\end{minted}

\begin{tikzpicture}[remember picture, overlay]

  \node[box, fit=(isdigs)(isdige)] (cisdig) {};
  \node[note, anchor=north west, text width=5cm]
    (noteIsdig) at ([xshift=-10cm, yshift=1.3cm] current page.east |- isdigs)
    {Retorna \texttt{True} si el carácter es un dígito del 0 al 9. Más directo que comparar con caracteres o rangos.};
  \draw[arr, dotted] (cisdig.east) to[bend right=5] (noteIsdig.west);

  \node[box, fit=(gets)(gete)] (cget) {};
  \node[note, anchor=north west, text width=5cm]
    (noteGet) at ([xshift=-10cm, yshift=-0.5cm] current page.east |- gets)
    {Si la clave aún no está en el diccionario, retorna \texttt{0} en vez de lanzar un error. Evita necesitar un \texttt{if} adicional para armar el histograma.};
  \draw[arr, dotted] (cget.west) to[bend right=5] (noteGet.west);

\end{tikzpicture}
\end{sectionbox}

Nota: Este ejercicio es un recorrido total. Ahora que saben usar \texttt{for}, esa también era una forma de resolverlo.

\end{document}
